% archivo: tabla_estacionariedad.tex
\begin{table}[htbp]
  \centering
  \begin{threeparttable}
    \caption{Contrastes de raíz unitaria sobre las series transformadas}
    \label{tab:estacionariedad}
    \begin{tabular}{l c c}
      \toprule
      Variable & ADF (p‑valor) & KPSS (p‑valor) \\
      \midrule
      ROA$_t$                          & 0.01 & 0.10 \\
      $\Delta L_t$ (Liquidez)          & 0.01 & 0.10 \\
      $\Delta T_t$ (Títulos)           & 0.01 & 0.10 \\
      $g^r_t$ (Remesas)                & 0.01 & 0.01 \\
      $\Delta i^p_t$ (Tasa pasiva)     & 0.01 & 0.10 \\
      $g^y_t$ (IMAE)                   & 0.01 & 0.10 \\
      \bottomrule
    \end{tabular}
    \begin{tablenotes}
      \item Nota: ADF: H0 = raíz unitaria; KPSS: H0 = estacionariedad. Valores $<0.05$ rechazan H0 al 5\%. Todas las series fueron desestacionalizadas (X‑13ARIMA‑SEATS) y diferenciadas. $g^r_t$ muestra evidencia mixta en KPSS, común en tasas financieras de alta frecuencia; se asume estacionariedad avalada por ADF.
    \end{tablenotes}
  \end{threeparttable}
\end{table}